% ======================================================================
%  31_flesh_hitchin.tex
%  骨血层之二：$\MH$ 作为二维靶空间；S-对偶 $=$ 镜像对称
%  主链 S3（§4 紧化）、S4（§5 的二维 TFT 与 S-对偶）；节点 N4（§4）、N5（§5）；miracle M2
%  底本：§4（印刷 pp.38--55 → p039--p056）；§5.3--5.6（印刷 pp.70--77 → p071--p078）
%  只写本文件；不改 preamble / main / 他人文件。
%  约定：〔补〕= 原文省略、由本节重建的推导，均带 \pB 或 \pC 与锚点。
% ======================================================================

\section{骨血层之二：Hitchin 模空间与 S-对偶即镜像对称}
\label{sec:flesh-hitchin}

\begin{kcard}{本节的位置：主链 S3--S4，节点 N4--N5，miracle M2}
\anchor{§4 开篇 p.38；§5.3 开篇 p.70；§5.6 收束 p.78}
\textbf{一句话。}把四维扭化 $\mathcal N=4$ 超对称规范理论放在 $M=\Sigma\times C$ 上，
并让 $C$ 远小于 $\Sigma$，低能有效理论就是 $\Sigma$ 上以 $\MH(G,C)$ 为靶空间的
二维超对称 sigma 模型；$\MH$ 是超凯勒流形，于是在它上面「取一个复结构做 B-模型、
取另一个复结构做 A-模型」有整整一个 $\PP^1$ 的自由度。四维的 S-对偶
$\tau\mapsto-1/\tau$ 在这套语言里变成：Hitchin 纤维化「基空间上的同构 $+$ 纤维上的
T-对偶」，也就是镜像对称。这是 \textbf{S3} 与 \textbf{S4}，合起来是 miracle \textbf{M2}。
\pA\anchor{src p.39}\pA\anchor{src p.70}

\textbf{输入与输出。}输入是 S2 的产物：一整个以 $t\in\PP^1$ 为参数的扭化拓扑场论族，
以及它的拓扑超荷 $Q$（$Q^2=0$，见骨血层之一）。输出是 S4 的产物：$\MH$ 上的膜范畴、
A/B 模型相对于三个复结构的分类（原文 Table 2），也就是 S5（线算子 $=$ 膜算子）
所需要的几何舞台。\pB\anchor{src p.56}

\textbf{为什么这一段必须慢读。}本节的两个结论都是一句话能说完的，正因如此极易被当成
「显然」跳过。原文 §4.1 用三页写下 $\MH$ 的超凯勒商结构，§4.2 只用一段话把三个复结构
与「Higgs 对 / 平坦 $\gC$-丛」对应起来，§5.4--5.5 又把最关键的论证放在「复结构 $I$ 的
B-模型的 chiral ring」里做，只在末尾补一句「我们要用的其实是 $J$ 的 B-模型与 $K$ 的
A-模型」。跳过任何一处，「S-对偶 $=$ 镜像对称」就只剩口号。
\end{kcard}

\subsection*{符号约定与新增符号}

本文件沿用 \texttt{02\_符号与风格卡} 的固定符号（$G$、$\Ldual$、$\gC$、$C$、$g$、
$\MH$、$\MC$、$\tau$、$t$、$\Psi$、$\varphi$、$A$、$F$、$E$、$Q$、$S$、$T$ 等）。
两点说明：原文 §4 的 $\varphi$ 与本文件 §2--§3 的「扭化 $1$-形式标量」是同一个场的
同一批分量，统一写 $\varphi$；原文 §4 的花体 $\mathcal M$（$C$ 上稳定 $G_{\CC}$-丛的
模空间）本文件记作 $\MC$，以免与 $\MH$ 混淆。下表只列新增符号，首次出现处再定义一次。

\begin{center}\small
\begin{tabular}{@{}lll@{}}
\toprule
\tabhead{符号} & \tabhead{含义} & \tabhead{出处}\\
\midrule
$\Sigma$ & 外部二维流形；$M=\Sigma\times C$ & \anchor{src p.38}\\
$W$ & 所有 $(A,\varphi)$ 构成的无穷维仿射空间（$C$ 上的位形空间） & \anchor{src p.40}\\
$I,J,K$ & $W$ 与 $\MH$ 上的三个复结构 & \anchor{src p.41}\\
$I_w$，$w\in\CC\cup\{\infty\}$ & $\PP^1_h$ 参数族的复结构，$I_0=I$，$I_{-\ii}=J$，$I_1=K$ & \anchor{src p.43}\\
$\omega_I,\omega_J,\omega_K$ & 与三个复结构相容的辛形式 & \anchor{src p.44}\\
$\Omega_I=\omega_J+\ii\omega_K$ & 复结构 $I$ 的全纯辛形式（循环定义 $\Omega_J,\Omega_K$） & \anchor{src p.44}\\
$\mu_I,\mu_J,\mu_K$ & 超凯勒商的三个矩映射 & \anchor{src p.45}\\
$\nu_I=\mu_J+\ii\mu_K$ & $G_{\CC}$ 作用的全纯矩映射 & \anchor{src p.49}\\
$B$，$\pi$ & Hitchin 纤维化的基（Casimir 空间）与映射 $\pi:\MH\to B$ & \anchor{src p.50}\\
$H_{P,\alpha}$ & 可积系统的 Hamiltonians & \anchor{src p.51}\\
$\mathcal F$，$\mathcal F_b$ & 一般纤维 $\pi^{-1}(b)$（复环面） & \anchor{src p.52}\\
$(w_+,w_-)$ & 二维 TFT 的两个复结构在 $\PP^1_h\times\PP^1_h$ 中的坐标 & \anchor{src p.60}\\
$\Xi$ & S-对偶在 Hitchin 基上诱导的分次环同构 $B\to{}^{L}\!B$ & \anchor{src p.74}\\
$\mathrm{J}\mathcal F$ & $\mathcal F$ 上秩一平坦 $U(1)$-丛的模空间（对偶环面） & \anchor{src p.76}\\
$(B,A,A)$ 等 & 膜相对于 $I,J,K$ 的 A/B 类型 & \anchor{src p.77}\\
\bottomrule
\end{tabular}
\end{center}

\subsection*{本节补写清单}

本节的重建集中在下面九处；正文逐处用〔补〕标出，并给锚点与重建依据。

\begin{center}\small
\begin{longtable}{@{}p{2.6cm}p{3.4cm}p{7.2cm}@{}}
\toprule
\tabhead{原文位置} & \tabhead{原文写法} & \tabhead{我们补了什么}\\
\midrule
\endhead
p.39（§4 开篇） & “one arrives at the conclusion obtained in [18]” & 四维场按 $C$ 的模式展开：KK 截断、零模空间 $W$、能量按 $\Sigma$ 导数阶数 $(p,q)$ 分解 \pB\\
p.39--40 & “taking the real and imaginary parts of $F$” & 复曲率恒等式 $F_{A+\ii\varphi}=(F-\varphi\wedge\varphi)+\ii D\varphi$；两种写法的 $F$ 不同名同物 \pB\\
p.45--47 & “a short computation shows” & 三个矩映射的显式计算；三个零点与 $(4.1)$ 三条方程的逐条对应 \pB\\
p.46--48 & “the metric and $B$-field are\dots” & 从 $(4.12)$--$(4.14)$ 到 $\dd s^2_g=\operatorname{Im}\tau\,\dd s^2$、$B=-(\operatorname{Re}\tau)\omega_I$ 的 $4\pi$ 系数 \pB\\
p.49--50 & “in complex structure $I$, the equations $(4.1)$ are the equations for a Higgs bundle” & $G$-商与 $G_{\CC}$-商两种描述在复结构 $I$ 下的等价 \pB\\
p.50--56 & “$\dim\mathcal F=\dim B$” & $\alpha$ 的次数补偿使被积式为 $(1,1)$-型；Poisson 对易；纤维紧 $\Rightarrow$ 环面 \pB\\
p.60--62 & “the general model is a generalized B-model” & $(J_+,J_-)$ 退化为 A/B 模型的判据 $w_+=w_-$ 与 $w_-=-1/w_+$ \pB\\
p.70--71 & $(5.35)$ 与 “$t$ is mapped to $t'=-t(\tau/|\tau|)$” & S-矩阵作用在 $(w_+,w_-)$ 与 $\Psi$ 上的显式计算 \pA\pB\\
p.75--77 & “the support of an A-brane\dots at least middle-dimensional” & zerobrane 论证中维数夹逼与 $\mathcal F\cong \mathrm{J}\tilde{\mathcal F}$ 的推导 \pB\\
\bottomrule
\end{longtable}
\end{center}

\subsection{导言：一条主线与四个先记住的结论}

本节按主链走：先重建 S3（四维 $\to$ 二维，\S\ref{sec:fh-from4d}），再讲 $\MH$ 的几何
（\S\ref{sec:fh-hk} 超凯勒结构、\S\ref{sec:fh-fib} 纤维化），最后讲 S4
（\S\ref{sec:fh-spec} 特殊化、\S\ref{sec:fh-s} S-对偶在纤维化上的作用、
\S\ref{sec:fh-brane} 膜分类与概念结论）。请按这个顺序读；先记住下面四个结论，
它们会在后文被逐条兑现。

\begin{keyidea}
\textbf{四个先记住的结论。}
\begin{enumerate}
\item 四维理论紧化到 $\Sigma$ 上，靶空间不是 $C$、不是 $M$、也不是某个 Calabi--Yau
流形，而是 $C$ 上 Hitchin 方程解的模空间 $\MH(G,C)$。\pA\anchor{src p.39}
\item $\MH$ 是超凯勒流形：它带三个复结构 $I,J,K$ 与三个辛形式，因而带一整族
$\PP^1$ 的复结构。二维拓扑场论只由一对复结构 $(J_+,J_-)$ 决定，所以同一个 $\MH$
可以承载很多不同的 A-模型与 B-模型。\pA\anchor{src p.58}
\item Hitchin 纤维化 $\pi:\MH\to B$ 把 $\MH$ 变成可积系统：基 $B$ 是向量空间，纤维是
环面，纤维在 $\omega_J,\omega_K$ 下是 Lagrangian。S-对偶在这上面就是「基上的同构
$+$ 纤维上的 T-对偶」，这正是 SYZ 的镜像对称图像。\pA\anchor{src p.71}
\item 于是 S-对偶 $=$ 镜像对称这句话的精确含义是：\emph{同一个}（或互为对偶的）超凯勒
流形上，一个复结构的 B-模型与另一个复结构的 A-模型之间的等价；不是两个 Calabi--Yau
流形之间的等价。\pA\anchor{src p.70}\pA\anchor{src p.77}
\end{enumerate}
\end{keyidea}

% ======================================================================
\subsection{从四维到二维：为什么靶空间是 $\MH$（§4 开篇，pp.~38--39）}
\label{sec:fh-from4d}

\subsubsection*{动机：不紧化，二维语言就无从谈起}

全书的目标 G0 是把几何 Langlands 对偶读成量子场论的一条定理。到 §3 为止四维一侧已经
准备好（S2：以 $t\in\PP^1$ 为参数的扭化拓扑场论族，拓扑超荷 $Q^2=0$）；但几何 Langlands
的语言是二维的——$C$ 上的层、Hecke 算子、平坦联络、$\Dmod$-模。要让「Wilson 线算子
$=$ 二维对象」有意义，必须先用紧化让四维理论\emph{长出}一个二维世界。

原文的选择是 $M=\Sigma\times C$：$C$ 是做几何 Langlands 的那条紧 Riemann 面（亏格
$g>1$，保证简单群 $G$ 的一般平坦联络不可约），$\Sigma$ 是另一个二维流形且显著大于
$C$。\pA\anchor{src p.38} 三个理由说明这不是任意的：扭化理论是拓扑的、对度量无依赖，
故 $M$ 可以取乘积；$C$ 必须是 Langlands 涉及的那条曲线（'t~Hooft 谱问题、Hecke
特征层、$D$-模都定义在它上面）；还需要一个外部流形 $\Sigma$ 承载二维场论与其膜——
膜的边界条件、线算子的插入都发生在 $\Sigma$ 上（这是 S5 的内容）。

尺寸层级是全部推导的技术前提：$C$ 上的 Kaluza--Klein 模式质量约 $1/R_C$，故
$E\ll1/R_C$ 时只剩零模；而从 $\Sigma$ 到某空间的缓变映射（变化尺度 $\gg R_C$）作用量
正比于小量 $(\partial_\Sigma\Phi)^2R_C^2$，$\Sigma$ 大时可以任意小。于是 $C$ 小保证只有
零模、$\Sigma$ 大保证零模能在 $\Sigma$ 上缓变。反过来若 $C\gg\Sigma$，一般的非常数映射
作用量很大，低能理论会变成靶空间为 $\MH(G,\Sigma)$ 的 sigma 模型。\pA\anchor{src p.39}

最后指出失败模式：取 $C$ 为一点则 $M=\Sigma$、靶空间退化为一点、无膜可谈；取 $\Sigma$
为一点则只得到 $C$ 上的量子力学。两个方向都必须存在：一个提供内部几何（$C$、Hitchin
方程），一个提供二维时空（$\Sigma$、映射）。

\subsubsection*{〔补〕推导一：按 $C$ 的模式展开}

原文在 p.39 用一句 “one arrives at the conclusion obtained in [18]” 把整段推导带过。
重建的依据是原文给出的两端：四维的平方和恒等式（§3.4，pp.~28--34）与 Hitchin 模空间
的定义（p.40）。\pB\anchor{src p.39}

\textbf{第一步（场的拆分与 KK 截断）。}玻色场为 $(A_\mu,\varphi_\mu,\sigma,\bar\sigma)$：
$\varphi$ 是扭化得到的 $\mathrm{ad}$ 值 $1$-形式，$\sigma,\bar\sigma$ 是 $0$-形式。
在 $M=\Sigma\times C$ 上拆指标：
\begin{equation}
A=A_\Sigma+A_C,\quad \varphi=\varphi_\Sigma+\varphi_C,\quad
A_\Sigma=A_m\dd x^m,\quad A_C=A_a\dd y^a .
\label{eq:split}
\end{equation}
把每个场按 $C$ 上（带联络的）Laplacian 的本征模式展开，模式 $\mathbf n$ 的四维质量平方
约为 $\lambda_{\mathbf n}/R_C^2$。因此 $E\ll1/R_C$ 时只剩零模，带非零 $C$-动量的模式
全部被积掉。这是标准 Kaluza--Klein 截断，合法性来自 $C$ 的小尺度，而不来自超对称。

\textbf{第二步（二维的动力学变量与超对称）。}零模中 $A_\Sigma,\varphi_\Sigma$ 是
$\Sigma$ 上的 $1$-形式，成为二维动力学场；而每一点 $x\in\Sigma$ 处的 $(A_C,\varphi_C)$
是 $C$ 上的一个规范场加一个 $\mathrm{ad}$ 值 $1$-形式。把全体这样的对记成
\begin{equation}
W=\bigl\{(A,\varphi)\ \big|\ A\ \text{为}\ C\ \text{上联络},\ 
\varphi\in\Omega^1(C,\mathrm{ad}E)\bigr\},
\label{eq:W}
\end{equation}
即原文 §4.1 的出发点。四维扭化理论保留原 $16$ 个超对称中的 $8$ 个，在 $\Sigma$ 上分成
两组各 $4$ 个（二维左右手征各一组），即二维 $(4,4)$ 超对称；二维的规范群是 $C\to G$
的映射群，而超多重态的空间正是 $W$。\pA\anchor{src p.46}

\textbf{第三步（能量按 $\Sigma$ 导数阶数分解）。}把四维作用量（§3.4 的平方和形式）按
$(p,q)$ 分解（$p$ 是 $\Sigma$ 指标数，$q$ 是 $C$ 指标数）：$(0,2)$ 型项不含 $\Sigma$
导数，扮演二维\emph{势能}；$(1,1)$ 型项含一个 $\Sigma$ 导数，给出 sigma 模型动能；
$(2,0)$ 型项是 $A_\Sigma$ 的动能，在 $C$ 上规范群被 Hitchin 解完全破缺处，$A_\Sigma$
获得质量、红外被积掉。\pA\anchor{src p.47}

\textbf{第四步（低能 $=$ 靶空间为 $\MH$ 的 sigma 模型）。}$(0,2)$ 型势能的零点逐点要求
$(A_C,\varphi_C)$ 是 $C$ 上 Hitchin 方程的解；模去 $C$ 上规范变换，全体解构成
$\MH(G,C)$。$(A_C,\varphi_C)$ 沿 $\Sigma$ 的缓慢变化就是一个映射
\begin{equation}
\Phi:\ \Sigma\longrightarrow \MH(G,C),\qquad x\longmapsto [A_C(x),\varphi_C(x)] .
\label{eq:Phi}
\end{equation}
于是低能有效理论正是以 $\MH(G,C)$ 为靶空间的二维超对称 sigma 模型；$\Sigma$ 尺度远大于
$C$ 保证了映射可以缓变、作用量很小。原文的结论句是：\pA\anchor{src p.39}

\begin{origquote}
“So, as long as we do not encounter fields on $C$ with continuous gauge symmetries, the effective physics at long distances on $\Sigma$ is given by a supersymmetric sigma-model in which the target space is $\MH(G,C)$.” \anchor{src p.39}
\end{origquote}

\textbf{补充说明。}两处省略原文也都提到：$\sigma,\bar\sigma$ 在零模层面必须满足
$D\sigma=[\varphi,\sigma]=[\sigma,\sigma]=0$，最简单的一类解取 $\sigma=0$；
$G=U(1)$ 或 $U(N)$ 时每个 Hitchin 解带未破缺的 $U(1)$ 规范对称，低能理论里除 sigma
模型还有二维 $U(1)$ 超对称规范理论，本文假设 $G$ 半单正是为避开它。\pA\anchor{src p.39}

\subsubsection*{〔补〕推导二：Hitchin 方程为什么有两种写法}

原文 §4 开篇说：$C$ 上的方程可以写成 $F=D{*}\varphi=0$，取 $F$ 的实部与虚部后
写成 $(4.1)$。这句话里藏着一个\textbf{符号陷阱}：两处的 $F$ 不是同一个对象。
第一种写法里的 $F$ 是\emph{复联络} $A+\ii\varphi$ 的曲率，第二种写法（$(4.1)$）里的
$F$ 是\emph{实联络} $A$ 的曲率。原文用「取实虚部」一句话过渡，读者若不留神，会以为
$(4.1)$ 只有两条方程，进而数错模空间的维数。下面把它算清楚。\pB\anchor{src p.39}

设 $A$ 是 $G$-联络，$\varphi$ 是 $\mathrm{ad}$ 值 $1$-形式，令
\begin{equation}
\mathcal A:=A+\ii\varphi .
\label{eq:complex-connection}
\end{equation}
则 $\mathcal A$ 是 $\gC$ 值联络，其曲率为
\begin{align}
F_{\mathcal A}
&=\dd\mathcal A+\mathcal A\wedge\mathcal A
=\dd(A+\ii\varphi)+(A+\ii\varphi)\wedge(A+\ii\varphi)\notag\\
&=\underbrace{\bigl(\dd A+A\wedge A-\varphi\wedge\varphi\bigr)}_{\textstyle F_A-\varphi\wedge\varphi}
+\ii\underbrace{\bigl(\dd\varphi+A\wedge\varphi+\varphi\wedge A\bigr)}_{\textstyle D_A\varphi}
=\bigl(F_A-\varphi\wedge\varphi\bigr)+\ii\,D_A\varphi ,
\label{eq:curvature-split}
\end{align}
其中 $D_A=\dd+[A,\cdot\,]$ 是伴随丛上的协变外微分。于是
\begin{equation}
F_{\mathcal A}=0
\quad\Longleftrightarrow\quad
F_A-\varphi\wedge\varphi=0\ \ \text{且}\ \ D_A\varphi=0 .
\label{eq:curvature-zero}
\end{equation}
在 $C$ 上，$F_A-\varphi\wedge\varphi$ 是 $\mathrm{ad}$ 值 $2$-形式，而 $\dim_\RR C=2$，
所以它只有一个 $(1,1)$ 分量：$F_A-\varphi\wedge\varphi=0$ 是一条 $\mathrm{ad}$ 值方程
（$\dim G$ 个实方程）。$D_A\varphi=0$ 是 $(1,1)$ 型的一条复方程（$\dim G$ 个实方程）。
再加上四维拓扑方程中独立剩下的
\begin{equation}
D_A{*}\varphi=0 ,
\label{eq:dstar}
\end{equation}
三条合起来就是 $(4.1)$：
\begin{equation}
F_A-\varphi\wedge\varphi=0,\qquad D_A\varphi=0,\qquad D_A{*}\varphi=0 .
\label{eq:hitchin}
\end{equation}
反过来，$(4.1)$ 也正好是 $\mathcal A=A+\ii\varphi$ 的平坦性加上 $D_A{*}\varphi=0$。
这就是为什么原文开头可以写 “$F=D{*}\varphi=0$”：那里 $F$ 是 $F_{\mathcal A}$，一旦
取实虚部，就得到 $(4.1)$ 的三条。\pB\anchor{src p.39}\pB\anchor{src p.29}

\begin{warn}
不要把 $(4.1)$ 记成「两条方程」。它的实方程数是
$\dim G+\dim G+\dim G=3\dim G$，对应的实维数是
$\dim_\RR\MH=4(g-1)\dim G$（见后）。少了任何一条，模空间的维数都不对。
\end{warn}

\subsubsection*{结论的作用：$\MH$ 这个靶空间被谁消费}

到此为止，S3 的产出可以精确写成：\emph{二维有效理论是映射 $\Phi:\Sigma\to\MH(G,C)$ 的
sigma 模型}。这句话的作用在后文四处兑现：它给 §5.1 提供舞台（扭化要在靶空间上选一对
复结构 $(J_+,J_-)$，$\MH$ 的超凯勒结构提供整族 $\PP^1_h$，否则这一步无从谈起）；给
§5.3--5.6 提供参数（二维的度规与 $B$-场由 $\tau$ 决定，而 $\tau$ 的 S-对偶正是四维的
S-对偶）；给 §6 提供膜（Wilson 线对应空间填充膜、't~Hooft 线对应余各向同性膜，都要
相对于 $\MH$ 的复结构来定义）；给 §11 提供 A-膜范畴（M5 的「典范余各向同性 A-膜
$\cong$ $\Dmod$-模」就陈述在 $\MH$ 上）。\pA\anchor{src p.58}\pA\anchor{src p.56}

\paragraph{排除错误直觉之一：$\MH$ 不是「紧化的内部空间」。}
许多人第一次读到「$M=\Sigma\times C$ 紧化后靶空间是 $\MH$」，会把 $\MH$ 想象成某种
Calabi--Yau 内空间，以为它是被紧化掉的那几个维度。这是错的：被紧化的是 $C$，$C$ 是
二维实流形（复曲线），不是 $\MH$；$\MH$ 是\emph{二维有效理论的靶空间}，是 sigma 模型
映射的值域，它的点代表「$C$ 上 Hitchin 方程的一个解模规范变换」。后果有三条：(i)
$\MH$ 一般非紧、带奇点（可约 Hitchin 解对应连续规范对称），不能当作紧 Calabi--Yau
来算；(ii) 它的复维数是 $2(g-1)\dim G$，随 $g$ 与 $\dim G$ 增长，不是固定的内空间
维数；(iii) 它的超凯勒结构由 $C$ 的复结构与 $G$ 决定，方向是「$C$ 给定，$\MH$ 才被
确定」，不是 $\MH$ 的模决定 $C$ 的模。\pB\anchor{src p.38}\pA\anchor{src p.39}

\subsubsection*{本节自检}
\begin{enumerate}
\item 为什么必须要求 $\Sigma$ 远大于 $C$？如果把两个尺度反过来，低能理论变成什么？
\item $(4.1)$ 有几条实方程？它们分别来自四维方程的哪一部分？
\item 在第六步里，$(A_C,\varphi_C)$ 沿 $\Sigma$ 的「缓慢变化」为什么不会破坏
$(0,2)$ 型势能的零点条件？
\end{enumerate}

% ======================================================================
\subsection{超凯勒结构：三个复结构、三个辛形式（§4.1--4.2，pp.~40--50）}
\label{sec:fh-hk}

\subsubsection*{动机：为什么靶空间必须超凯勒}

\S\ref{sec:fh-from4d} 的结论是「靶空间 $=\MH$」。但二维拓扑场论要区分 A-模型与 B-模型，
靠的是靶空间上的复结构或辛结构：B-模型选一个复结构，A-模型选一个辛结构（等价地，
选一个复结构并取其 Kähler 形式；原文特别强调 A-模型真正依赖的是辛形式而不是复结构本身）。
如果 $\MH$ 只有一个复结构，那么同一个靶空间上就只有一种 B-模型与一种 A-模型，
「S-对偶交换 B 与 A」就失去了几何前提。M2 要成立，$\MH$ 必须带\emph{一整族}复结构。

超凯勒流形正好提供这一族：若 $X$ 带三个复结构 $I,J,K$（满足四元数关系）与一个在
三者下都为 Kähler 的度量，则任意单位向量 $(a,b,c)$ 给出新的复结构
$aI+bJ+cK$，全体构成一个 $\PP^1$；每一对点 $(J_+,J_-)$ 给出一个二维 $(2,2)$ 超对称
子代数。原文的策略是证明 $\MH$ 是超凯勒商 $W///G$，从而自动超凯勒。本小节按原文
§4.1 的顺序把结构写出来。

\subsubsection*{$W$ 上的平坦超凯勒结构}

先写度量。$W$ 是仿射空间（(eq:W)），其切空间在每一点可等同于变分 $(\delta A,\delta\varphi)$
的空间。原文的度量是\pA\anchor{src p.41}
\begin{equation}
\dd s^2=-\frac{1}{4\pi}\int_C|\dd^2z|\,\Tr\Bigl(
\delta A_z\otimes\delta A_{\bar z}+\delta A_{\bar z}\otimes\delta A_z
+\delta\varphi_z\otimes\delta\varphi_{\bar z}+\delta\varphi_{\bar z}\otimes\delta\varphi_z
\Bigr),\qquad |\dd^2z|:=\ii\,\dd z\wedge\dd\bar z .
\label{eq:flat-metric}
\end{equation}
它是正定的（$\Tr$ 在反 Hermitian 元素上为负，前面的负号把它翻正），并且只依赖 $C$ 的
复结构与 $G$ 的 Killing 形式，不依赖 $C$ 上的坐标选取。

三个复结构不必用分量矩阵写死；用\emph{全纯函数}来定义更清楚，而且这正是原文
§4.2 用来识别它们的判据。$W$ 上的复结构由它把哪些 $(1,0)$-形式定为全纯决定：\pA\anchor{src p.42}
\begin{itemize}
\item $I$：$A_z$ 与 $\varphi_z$ 全纯（等价地，$\delta A_z,\delta\varphi_z$ 是 $(1,0)$
  形式，$\delta A_{\bar z},\delta\varphi_{\bar z}$ 是 $(0,1)$ 形式）。$I$ 保持 $C$ 的
  $z/\bar z$ 分块，在 $z$ 块上作用如 $\ii$、在 $\bar z$ 块上如 $-\ii$。
\item $J$：$A+\ii\varphi$ 的全部分量全纯。等价地
  \(J^t\delta A_z=-\delta\varphi_{\bar z}\)，\(J^t\delta\varphi_z=\delta A_{\bar z}\)，
  于是 \(J^t(\delta A+\ii\delta\varphi)=\ii(\delta A+\ii\delta\varphi)\)。
  注意 $J$ \emph{交换} $z$ 与 $\bar z$ 两个分块。
\item $K$：$A_z-\varphi_z$ 与 $A_{\bar z}+\varphi_{\bar z}$ 全纯。
\end{itemize}
三条可以直接从原文 $(4.3)$--$(4.5)$ 的分量形式读出（底本的文本层丢了共轭号，
我们按 p.43 明确给出的三组全纯函数重新写成上面这种不依赖符号约定的形式）。
它们满足四元数关系
\begin{equation}
IJ=-JI=K,\qquad JK=-KJ=I,\qquad KI=-IK=J ,
\label{eq:quat}
\end{equation}
这正是「三个复结构而不是三个随便的复结构」的含义。\pA\anchor{src p.42}

超凯勒的第二个条件是度量为三个复结构共同的 Hermitian 度量。以 $J$ 为例，把
(eq:flat-metric) 的每一项用 $J^t$ 作用，并用 $J^t$ 交换两个分块这一事实：
\begin{align}
J^t\otimes J^t\bigl(\dd s^2\bigr)
&=-\frac{1}{4\pi}\int_C|\dd^2z|\,\Tr\Bigl(
\delta\varphi_{\bar z}\otimes\delta\varphi_z+\delta\varphi_z\otimes\delta\varphi_{\bar z}
+\delta A_{\bar z}\otimes\delta A_z+\delta A_z\otimes\delta A_{\bar z}\Bigr)\notag\\
&=\dd s^2 .
\label{eq:J-hermitian}
\end{align}
对 $I$ 与 $K$ 同样成立（$I$ 在 $z$ 块上给 $\ii$、在 $\bar z$ 块上给 $-\ii$，两个因子的
相位相乘为 $1$；$K$ 把 $\delta A_z$ 送到 $-\ii\delta\varphi_z$、把 $\delta A_{\bar z}$
送到 $\ii\delta\varphi_{\bar z}$，同样配对）。因此 $W$ 是一个\emph{平坦}超凯勒流形：
仿射空间、平度量、三个平移不变且两两正交的复结构。平坦性使得后面的所有验证都是
线性代数，不需要曲率计算。\pA\anchor{src p.42}

\begin{keyidea}
\textbf{为什么「平坦」很重要。}超凯勒商的一般理论只需要 $X$ 超凯勒、$H$ 以三个矩映射
作用；而本例中 $X=W$ 是仿射空间，三个复结构是常系数线性变换，度量是常系数二次型。
于是「$W$ 超凯勒」这件事可以逐项手算（如 \eqref{eq:J-hermitian}），
不必引用任何存在性定理。这是原文 §4.1 敢于只用三页讲完 $\MH$ 超凯勒的原因。
\end{keyidea}

\subsubsection*{$\PP^1_h$：一整族复结构}

$I,J,K$ 的任意单位线性组合都是复结构：
\begin{equation}
I_w=\frac{(1-|w|^2)\,I+\ii(w-\bar w)\,J+(w+\bar w)\,K}{1+|w|^2},
\qquad w\in\CC\cup\{\infty\}=:\PP^1_h .
\label{eq:Iw}
\end{equation}
容易核对：$w=0$ 给 $I$，$w=\infty$ 给 $-I$，$w=-\ii$ 给 $J$，$w=1$ 给 $K$，
$w=\ii$ 给 $-J$，$w=-1$ 给 $-K$；一般的 $w$ 对应 $\PP^1_h$ 上单位球面的一点
$(a,b,c)$，$a^2+b^2+c^2=1$。\pA\anchor{src p.43}

一个在后文反复使用的恒等式是\pB\anchor{src p.43}
\begin{equation}
I_{-1/w}=-I_w .
\label{eq:opposite}
\end{equation}
证明只需把 $w\mapsto-1/w$ 代入 \eqref{eq:Iw} 并通分：
\begin{align}
I_{-1/w}
&=\frac{(1-|w|^{-2})I+\ii(-1/w+1/\bar w)J+(-1/w-1/\bar w)K}{1+|w|^{-2}}\notag\\
&=\frac{(|w|^2-1)I+\ii(\bar w-w)J-(\bar w+w)K}{1+|w|^2}
=-\frac{(1-|w|^2)I+\ii(w-\bar w)J+(w+\bar w)K}{1+|w|^2}=-I_w .
\label{eq:opposite-proof}
\end{align}
\eqref{eq:opposite} 是 §\ref{sec:fh-spec} 中「A-模型判据」的全部代数内容：
两个复结构相反 $\iff$ $w_-=-1/w_+$。

另外，$W$ 上还有一个 $U(1)$ 作用：固定 $A$，把 $\varphi_z\mapsto\lambda\varphi_z$、
$\varphi_{\bar z}\mapsto\lambda^{-1}\varphi_{\bar z}$，$|\lambda|=1$。它保持度量与三个
辛形式，在 $\PP^1_h$ 上作用为 $I_w\mapsto I_{\lambda^{-1}w}$，固定 $w=0$ 与 $w=\infty$。
这个 $U(1)$ 不是四维理论的对称群，但 $U(1)$ 的 $-1$ 元（即 $\varphi\mapsto-\varphi$）
来自四维扭化理论：它对应 $SU(4)_R$ 的中心元，在 $\mathbf 4$ 上为 $\ii$、在
$\mathbf 6=\wedge^2\mathbf 4$ 上为 $-1$。$\varphi\mapsto-\varphi$ 在复结构 $I$ 下全纯、
在 $J,K$ 下反全纯，其不动点集里 $\varphi=0$ 那一支给出稳定丛模空间
$\mathcal M\hookrightarrow\MH$ 的嵌入。\pA\anchor{src p.43}

\subsubsection*{三个辛形式与全纯辛形式}

每个复结构与度量给出一个辛形式 $\omega_I(\cdot,\cdot)=\dd s^2(I\cdot,\cdot)$，循环定义
$\omega_J,\omega_K$。原文的显式公式是\pA\anchor{src p.44}
\begin{align}
\omega_I&=-\frac{\ii}{2\pi}\int_C|\dd^2z|\,\Tr\bigl(\delta A_z\wedge\delta A_{\bar z}-\delta\varphi_z\wedge\delta\varphi_{\bar z}\bigr)
=-\frac{1}{4\pi}\int_C\Tr\bigl(\delta A\wedge\delta A-\delta\varphi\wedge\delta\varphi\bigr),\label{eq:omegaI}\\
\omega_J&=\frac{1}{2\pi}\int_C|\dd^2z|\,\Tr\bigl(\delta\varphi_z\wedge\delta A_{\bar z}+\delta\varphi_{\bar z}\wedge\delta A_z\bigr),\label{eq:omegaJ}\\
\omega_K&=\frac{\ii}{2\pi}\int_C|\dd^2z|\,\Tr\bigl(\delta\varphi_z\wedge\delta A_{\bar z}-\delta\varphi_{\bar z}\wedge\delta A_z\bigr)
=\frac{1}{2\pi}\int_C\Tr\,\delta\varphi\wedge\delta A .\label{eq:omegaK}
\end{align}
三条性质值得记住。第一，$\omega_J$ 依赖 $C$ 的复结构（它把 $\varphi_z$ 与 $A_{\bar z}$
配对），而 $\omega_I,\omega_K$ 不依赖——这与「$J$ 是特殊复结构」是同一件事的两面。
第二，$\omega_J,\omega_K$ 在 $W$ 上恰当（exact），而 $\omega_I$ 上同调非平凡；其
上同调类与行列式线丛的曲率成正比，原文用 (4.10) 把它写成 $\omega_I=\omega_I'+\delta\lambda_I$。
第三，定义复化组合
\begin{equation}
\Omega_I:=\omega_J+\ii\omega_K,\qquad
\Omega_J:=\omega_K+\ii\omega_I,\qquad
\Omega_K:=\omega_I+\ii\omega_J ,
\label{eq:Omega}
\end{equation}
则 $\Omega_I$ 是复结构 $I$ 下的 $(2,0)$ 形式且全纯闭，称为 $\MH$ 在 $I$ 下的\emph{全纯辛形式}。
它对后文的纤维化与膜分类都是关键：Hitchin 纤维的 Lagrangian 性、$(B,A,A)$-膜的构造，
都要用它来陈述。\pA\anchor{src p.44}

\subsubsection*{矩映射与超凯勒商：$\MH=W///G$}

规范变换群以 $\delta A=-D\epsilon$、$\delta\varphi=[\epsilon,\varphi]$ 作用在 $W$ 上
（$\epsilon$ 是 $\mathrm{ad}$ 值 $0$-形式）。它保持度量与三个辛形式，因此若三个矩映射
都存在，就得到超凯勒作用。矩映射由
$\iota_{V(\epsilon)}\omega=\delta\mu$ 定义（$V(\epsilon)$ 为 $\epsilon$ 生成的向量场）。
把 \eqref{eq:omegaI} 代入直接计算，我们得到\pB\anchor{src p.45}
\begin{align}
\iota_{V(\epsilon)}\omega_I
&=\frac{1}{2\pi}\int_C\Tr\bigl(D\epsilon\wedge\delta A+[\epsilon,\varphi]\wedge\delta\varphi\bigr)
=\frac{1}{2\pi}\int_C\Tr\bigl(-\epsilon\,\delta F+\epsilon\,[\varphi,\delta\varphi]\bigr)\notag\\
&=-\delta\left[\frac{1}{2\pi}\int_C\Tr\epsilon\,\bigl(F-\varphi\wedge\varphi\bigr)\right],
\label{eq:muI-derivation}
\end{align}
其中第二步用了分部积分（伴随丛上 $D$ 的伴随为 $-D$，边界项略去）与 $\Tr$ 的循环性，
第三步用了 $\delta(F-\varphi\wedge\varphi)=\delta F-2\,\delta\varphi\wedge\varphi$。于是
\begin{equation}
\mu_I=-\frac{1}{2\pi}\int_C\Tr\epsilon\,(F-\varphi\wedge\varphi),
\label{eq:muI}
\end{equation}
与原文 $(4.11)$ 一致。对 $\omega_J,\omega_K$ 做同样的计算给出另外两个矩映射；把它们
合并成全纯矩映射最方便：\pB\anchor{src p.46}
\begin{equation}
\nu_I:=\mu_J+\ii\mu_K=-\frac{1}{\pi}\int_C|\dd^2z|\,\Tr\epsilon\,D_{\bar z}\varphi_z .
\label{eq:nuI}
\end{equation}
（底本的文本层把 \eqref{eq:nuI} 中的共轭号丢了；按 $z\leftrightarrow\bar z$ 的对称性
与 p.49 的几何解释——$\nu_I=0$ 当且仅当 $\varphi_z\,\dd z$ 是全纯截面——可以确定
这里应当是 $D_{\bar z}$。下文只用这个几何解释，不用具体的分量式。）

三个矩映射同时为零的条件现在可以逐条读出：
\begin{itemize}
\item $\mu_I=0$ $\iff$ $F-\varphi\wedge\varphi=0$（(eq:hitchin) 的第一条）；
\item $\nu_I=0$ $\iff$ $D_{\bar z}\varphi_z=0$，即 $\varphi=\varphi_z\,\dd z$ 是
  $K_C\otimes\mathrm{ad}E$ 的全纯截面（这就是「Higgs 对」的条件）；
\item $\mu_J=\mu_K=0$ 的实部与虚部合起来给出 $D\varphi=0$ 与 $D{*}\varphi=0$
  （(eq:hitchin) 的后两条）。
\end{itemize}
所以三个矩映射的零点正是 $(4.1)$ 的三条 Hitchin 方程。\pB\anchor{src p.46}

超凯勒商的一般定理现在可以用了：若超凯勒流形 $X$ 上李群 $H$ 以三个矩映射
$\vec\mu=(\mu_I,\mu_J,\mu_K)$ 作用，则 $\vec\mu^{-1}(0)/H$ 是超凯勒流形，称为
\emph{超凯勒商} $X///H$。取 $X=W$、$H=G$，即得
\begin{equation}
\MH(G,C)=W///G=\bigl\{\mu_I=\mu_J=\mu_K=0\bigr\}\big/G .
\label{eq:HKquotient}
\end{equation}
这就是 $\MH$ 超凯勒的来源。\pA\anchor{src p.46}

\begin{remark}[两种商描述]
超凯勒商有两种等价写法：先解 $\mu_I=0$（实矩映射）再除以 $G$，或者把 $G_{\CC}$ 的
复化作用作用在复化后的 $W$ 上、解全纯矩映射 $\nu_I=0$ 再除以 $G_{\CC}$。在复结构 $I$
下后一种给出 Higgs 对的模空间 $\{D_{\bar z}\varphi_z=0\}/G_{\CC}$，原文 p.49--50 明确
以它作为「复结构 $I$ 就是 Higgs 丛的复结构」的理由；这也是 §\ref{sec:fh-brane} 膜分类
的语言基础。\pA\anchor{src p.49}
\end{remark}

维数核验：$\dim_\RR\MH=4(g-1)\dim G$，即 $\dim_\CC\MH=2(g-1)\dim G$，与原文
p.40 的维数一致。这个数也可以从 Hitchin 纤维化独立算出（下一小节），两者的吻合是
后面所有维数论证的基准。\pA\anchor{src p.39}

\subsubsection*{为什么二维理论里「只有 $I$ 显式出现」}

这是全节最容易被跳过、却决定后续一切的地方。原文 §4 从头到尾几乎只用 $I$：
§4.2 把 $\MH$ 描述成 Higgs 对的模空间（复结构 $I$），§4.3 的 Hitchin 纤维化对 $I$
全纯，$B$-场公式 $B=-(\operatorname{Re}\tau)\omega_I$ 也只用到 $\omega_I$。读者容易
由此以为「物理上重要的复结构是 $I$」。恰恰相反。正确的读法是三层。

\textbf{第一层（物理模型在 $J$ 与 $K$）。}四维族在 §5.3 的特殊化给出：$\Psi=\infty$
时二维模型是 $J$ 中的 B-模型，$\Psi=0$ 时是 $K$ 中的 A-模型（见 §\ref{sec:fh-spec}）。
$I$ 中的 B-模型与 A-模型虽然也在族里，但它们是 S-对偶的不动点（$I_B,I_A$ 在 Table 2
中各自对自己），不参与「B 换成 A」的交换。\pA\anchor{src p.70}

\textbf{第二层（$I$ 是计算用的复结构）。}$I$ 是唯一一个使 $\MH$ 变成\emph{代数}对象
的复结构：在 $I$ 下 $\MH$ 是 Higgs 丛的模空间，Hitchin 纤维化是代数的，纤维是环面，
chiral ring 就是基 $B$ 上的全纯函数。所有这些让「S-对偶在基上如何作用」可以算出来。
原文在 §5.4 末尾对此有一句坦白的说明：\pA\anchor{src p.73}
\begin{origquote}
“We made the argument in the B-model of complex structure $I$ even though we will apply the results to the B-model in complex structure $J$ and the A-model in complex structure $K$.” \anchor{src p.73}
\end{origquote}
换言之，$I$ 是\emph{证明工具}，不是物理结论。这正是读者最容易误解的地方。

\textbf{第三层（为什么 $I$ 是「显式」的而 $J,K$ 不是）。}$I$ 保持 $C$ 的 $z/\bar z$
分块（$K$ 也保持，但配对方式不同），$J$ 交换两个分块、因此不依赖 $C$ 的复结构
（p.43 特别指出）。物理上，$B$-场沿 $\omega_I$ 是因为 $\omega_I$ 在 $I$ 下是 $(1,1)$
型；在 $J$ 的 B-模型中 $\omega_I$ 的 $(2,0),(0,2)$ 分量被 $Q$-恰当项吸收，只有
$(1,1)$ 分量起作用（p.70）。所以「公式里只出现 $I$」反映的是 B-场的 Hodge 型，不是
物理模型的选择。\pB\anchor{src p.70}

\paragraph{排除错误直觉之二：三个复结构不是三个理论，A-模型也不「用复结构」。}
$\MH$ 上的三个复结构 $I,J,K$ 不是三个不同的二维理论，而是\emph{同一个}超凯勒流形上
三个可供选择的复结构；一个二维 TFT 由一对 $(J_+,J_-)$ 决定，B-模型对应 $J_+=J_-$，
A-模型对应 $J_-=-J_+$（见 §\ref{sec:fh-spec}）。所以「$J$ 中的 B-模型」与「$K$ 中的
A-模型」是同一个 $\MH$ 上的两个不同 TFT，而不是两个不同的流形。另外，A-模型真正依赖
的是辛形式 $\omega$ 而不是复结构本身；说「$K$ 中的 A-模型」只是简称，准确的说法是
「以 $\omega_K$ 为辛形式的 A-模型」。这一点在膜的分类里会造成实际差别：
A-膜的支撑只需对 $\omega_K$ 是 Lagrangian（或余各向同性），不必是 $K$-全纯的。\pA\anchor{src p.59}

\subsubsection*{图：超凯勒的三个复结构与 $\PP^1_h$}

\begin{figure}[htbp]
\centering
\begin{tikzpicture}[scale=0.74,>={Stealth[length=2mm]}]
  % 球面
  \fill[blue!4] (0,0) circle (2);
  \draw[thick] (0,0) circle (2);
  % 赤道：正面实线、背面虚线
  \draw[gray!70] (-2,0) arc (180:360:2 and 0.62);
  \draw[gray!70,dashed] (2,0) arc (0:180:2 and 0.62);
  % 南北轴
  \draw[gray,dashed] (0,-2)--(0,2);
  % 六个特殊点：w=0,∞ 在南北极；w=±i,±1 在赤道上
  \fill (0,2) circle (1.7pt) node[above=1pt] {$I$};
  \fill (0,-2) circle (1.7pt) node[below=1pt] {$-I$};
  \fill (2,0) circle (1.7pt) node[right=2pt] {$K$};
  \fill (-2,0) circle (1.7pt) node[left=2pt] {$-K$};
  \fill (0,0.62) circle (1.7pt) node[right=3pt] {$J$};
  \fill (0,-0.62) circle (1.7pt) node[right=3pt] {$-J$};
  % U(1) 轨道：绕 I 轴的小圆
  \draw[->,thick,kwblue] (0.85,1.15) arc (0:330:0.85 and 0.24);
  \node[kwblue,font=\small] at (1.55,1.42) {$U(1)_R$};
  % w 标注
  \node[gray,font=\scriptsize,align=left] at (-1.15,2.35) {$w=0$};
  \node[gray,font=\scriptsize,align=left] at (-1.05,-2.35) {$w=\infty$};
  \node[gray,font=\scriptsize] at (2.45,0.28) {$w=1$};
  \node[gray,font=\scriptsize] at (1.45,0.75) {$w=-\ii$};
\end{tikzpicture}
\caption{超凯勒流形的复结构族。球面是
$\{aI+bJ+cK\mid a^2+b^2+c^2=1\}\cong\PP^1_h$，三个正交的复结构 $I,J,K$ 是球面上
三个互相垂直的点；$U(1)_R$ 作用绕 $I$ 轴旋转，固定 $I$ 与 $-I$。按 \eqref{eq:Iw}，
$w=0,-\ii,1$ 分别给出 $I,J,K$。4d 族中「$\Psi=\infty$ 的 B-模型」落在 $J$（$w_+=w_-=-\ii$），
「$\Psi=0$ 的 A-模型」落在赤道大圆上的相反点对 $(w_+,w_-)=(-t,t^{-1})$，$t$ 实。}
\label{fig:hk-sphere}
\end{figure}

\subsubsection*{本节自检}
\begin{enumerate}
\item 为什么「$\MH$ 超凯勒」是 M2（S-对偶 $=$ 镜像对称）的必要条件？如果 $\MH$ 只有单个
复结构，S-对偶交换的是什么？
\item 用 \eqref{eq:Iw} 验证 $w=\ii$ 给出 $-J$、$w=-1$ 给出 $-K$。
\item $J^t\delta A_z=-\delta\varphi_{\bar z}$ 为什么等价于「$A+\ii\varphi$ 全纯」？
这个式子说明了 $J$ 与 $I$ 在对待 $C$ 的复结构上有什么不同？
\end{enumerate}

% ======================================================================
\subsection{$\MH$ 上的 Hitchin 纤维化：为什么它是可积系统（§4.3，pp.~50--56）}
\label{sec:fh-fib}

\subsubsection*{动机：把 $\MH$ 变成可计算的对象}

到 §\ref{sec:fh-hk} 为止，$\MH$ 只是一个「超凯勒商」：定义清楚，但不可计算。
它是非紧的、带奇点的、维数随 $g$ 增长的复流形；直接在上面做膜论或做 S-对偶，
无从下手。原文 §4.3 的任务是给 $\MH$ 一个\emph{可计算的结构}：一个到向量空间的
纤维化，其纤维是环面。这个结构同时是三个东西：
\begin{enumerate}
\item 一个\emph{代数完全可积系统}（Hitchin 1987）：基空间的线性函数就是互相对易的
  Hamiltonians，纤维是其等值面；
\item 一个\emph{特殊 Lagrangian 环面纤维化}：纤维在 $\omega_J,\omega_K$ 下是 Lagrangian，
  在复结构 $I$ 下是全纯子流形——这正是 SYZ 镜像对称所需的几何（§\ref{sec:fh-s}）；
\item 一个\emph{分级环的载体}：$\MH$ 在复结构 $I$ 下的全纯函数就是基 $B$ 上的全纯函数，
  因而 B-模型的 chiral ring 可以显式算出来（§\ref{sec:fh-s}）。
\end{enumerate}
第二、三条是本节全部技术工作的目的。原文在 §5.4 开头一句话点明：
「我们的目标是证明 S-对偶在 Hitchin 纤维化的基上经典地作用、在纤维上以 T-对偶作用」
（原句见 §\ref{sec:fh-s}）。没有 §4.3 的纤维化，这句话无法陈述。

\subsubsection*{构造：Casimir 不变量给出映射 $\pi:\MH\to B$}

设 $G$ 的秩为 $r$，$\mathfrak g$ 的基本不变多项式为 $P_1,\dots,P_r$，次数为
$d_1,\dots,d_r$（例如 $G=SU(N)$ 时 $d_i=i+1$，$i=1,\dots,N-1$）。Hitchin 映射定义为
\begin{equation}
\pi:\ \MH(G,C)\longrightarrow B:=\bigoplus_{i=1}^{r}H^0\bigl(C,K_C^{\,d_i}\bigr),
\qquad (E,\varphi)\longmapsto\bigl(P_1(\varphi),\dots,P_r(\varphi)\bigr).
\label{eq:hitchin-map}
\end{equation}
这里 $P_i(\varphi)$ 是把 $\varphi\in H^0(K_C\otimes\mathrm{ad}E)$ 代入不变多项式得到的
截面：$\varphi$ 的每个分量在 $K_C$ 中，$n$ 次多项式给 $K_C^{\,n}$ 中的截面。
在复结构 $I$ 下 $\pi$ 是全纯的；在复结构 $J$（平坦 $\gC$-丛）下 $\pi$ 有另一个解释
（monodromy 的特征值），但本节只用 $I$。\pA\anchor{src p.50}

维数计算是理解整个纤维化的关键，值得逐步写清。由 Riemann--Roch，
$\dim H^0(C,K_C^{\,n})=(2n-1)(g-1)$（$n\ge2$；基本不变量的次数 $d_i\ge2$，
所以 $n=1$ 的情形不会出现，无需担心 $K_C$ 的一次幂）。于是
\begin{equation}
\dim_\CC B=\sum_{i=1}^{r}(2d_i-1)(g-1)=(g-1)\sum_{i=1}^{r}(2d_i-1)=(g-1)\dim G ,
\label{eq:dimB}
\end{equation}
最后一步用了简单 Lie 代数的经典恒等式 $\sum_i(2d_i-1)=\dim\mathfrak g$。
而由 §\ref{sec:fh-hk}，$\dim_\CC\MH=2(g-1)\dim G$，所以
\begin{equation}
\dim_\CC B=\tfrac12\dim_\CC\MH .
\label{eq:halfdim}
\end{equation}
一个复维数一半的靶空间——这正是完全可积系统的维数特征。\pA\anchor{src p.50}

\subsubsection*{基的线性函数就是互相对易的 Hamiltonians}

原文给出 Hamiltonians 的显式形式：对每个不变多项式 $P$（次数 $n$）与
$\alpha\in H^1(C,K_C^{\,1-n})$，令\pA\anchor{src p.51}
\begin{equation}
H_{P,\alpha}=\int_C\alpha\,P(\varphi).
\label{eq:H}
\end{equation}
次数配平需要看清：$P(\varphi)$ 是 $K_C^{\,n}$ 的截面，而 $\alpha$ 按 Serre 对偶
是 $K_C^{\,1-n}$ 值的 $(0,1)$-形式；两者相乘得到 $(1,1)$-形式，才能在 $C$ 上积分。
而
\begin{equation}
\dim H^1\bigl(C,K_C^{\,1-n}\bigr)=\dim H^0\bigl(C,K_C^{\,n}\bigr)=(2n-1)(g-1),
\label{eq:serre}
\end{equation}
即 $\alpha$ 的空间恰是 $H^0(C,K_C^{\,n})^*$。换句话说，$H_{P,\alpha}$ 作为
$\bigl(P(\varphi)\bigr)$ 的线性函数，正是基 $B$ 上的一个\emph{线性坐标}。所以：
\begin{center}
\emph{基 $B$ 的线性函数 $\Longleftrightarrow$ 互相对易的 Hamiltonians。}
\end{center}
这解释了 \eqref{eq:halfdim} 的另一层含义：可积系统所需的 $\tfrac12\dim\MH$ 个独立
Hamiltonians，正好是 $B$ 的维数个线性函数。\pB\anchor{src p.51}

为什么它们互相对易？原文的论证分两步。第一步，在 $T^*\mathcal M\subset\MH$（$\varphi=0$
截面的余切方向的稠密开集，即「Hitchin 第二纤维化」给出的图卡）上，全纯辛形式是
$\Omega_I=\int_C\Tr\delta\varphi_z\wedge\delta A_z$，所以 $\varphi_z$ 与 $A_z$ 是共轭
变量。而 $H_{P,\alpha}$ 只依赖 $\varphi$（不依赖 $A$），因此任意两个这样的函数在
$\Omega_I$ 下的 Poisson 括号为零。第二步，Poisson 括号在超凯勒商下良定义
（商掉的是 Hamiltonian 作用），所以 $T^*\mathcal M$ 上的对易性延拓到整个 $\MH$。
原文把它总结为「$\MH$ 上存在 $\tfrac12\dim\MH$ 个对易的全纯函数，它们生成 $\MH$ 的
全纯函数环」。\pB\anchor{src p.52}

\subsubsection*{纤维：紧、全纯、Lagrangian、环面}

现在看映射 $\pi$ 的纤维 $\mathcal F_b=\pi^{-1}(b)$。原文给出的性质如下：\pA\anchor{src p.53}
\begin{itemize}
\item $\mathcal F_b$ 是紧复子流形（在复结构 $I$ 下全纯），复维数为
  $\dim_\CC\MH-\dim_\CC B=(g-1)\dim G$；
\item $\mathcal F_b$ 是复环面的 torsor（一个 torus 的平移），因此带交换群结构；
\item $\mathcal F_b$ 在 $\omega_J$ 与 $\omega_K$ 下是 Lagrangian 子流形
  （middle-dimensional，且两个辛形式在纤维上拉回为零）。
\end{itemize}
第三条是后文 SYZ 论证的关键。为什么会 Lagrangian？因为纤维由对易 Hamiltonians 的
等值面 $H_k=h_k$ 定义，是 middle-dimensional 的；而对易 Hamiltonians 的等值面在
Liouville--Arnold 意义下是 Lagrangian 的（原文 p.56 把这个推理一句带过，我们在这里
标为 \pB）。把它翻译成超凯勒语言就是：纤维在复结构 $I$ 下全纯（所以是复子流形），
又在 $\omega_J,\omega_K$ 下迷向（Lagrangian），这正是 SYZ 纤维化的要求。
\pB\anchor{src p.56}

\begin{keyidea}
\textbf{一句话记住 Hitchin 纤维化。}基 $B$ 是「可测量」的部分（Casimir 不变量，
经典可观测量的值），纤维 $\mathcal F_b$ 是「不可测量」的部分（环面，只能量子地感知）。
这句话在 §\ref{sec:fh-s} 会被原文用两个论证独立确认：chiral ring 论证（基上的全纯函数
就是全部局域算符）与 zerobrane 论证（点膜在纤维上移动，镜像是纤维上的 T-对偶）。
\end{keyidea}

\subsubsection*{附：Hitchin 第二纤维化与 $\mathcal M\hookrightarrow\MH$}

除了 $\pi:\MH\to B$，原文还用到第二个纤维化：忘掉 $\varphi$，把 $(E,\varphi)$ 送到
$E\in\MC$。它只在 $\varphi=0$ 附近定义良好，纤维是 $T^*\MC$ 的平移，而 $T^*\MC$ 是
$\MH$ 的稠密开集。两个纤维化横截：第一纤维化的纤维（环面）与第二纤维化的纤维
（仿射空间）互补。第二纤维化在推导 $B$-场的 $4\pi$ 系数与 §\ref{sec:fh-brane} 的
$(B,A,A)$-膜例子中都要用到；$\MC$ 在复结构 $I$ 下全纯、在 $\omega_J,\omega_K$ 下
Lagrangian，是「special Lagrangian」在超凯勒情形的类比物。\pA\anchor{src p.54}

\subsubsection*{图：纤维化与纤维-底空间}

\begin{figure}[htbp]
\centering
\begin{tikzpicture}[scale=0.72,>={Stealth[length=2mm]}]
  % 总空间外框
  \draw[fill=blue!3,draw=blue!45,rounded corners=7pt] (-4.3,0.60) rectangle (4.3,3.50);
  \node[font=\small,kwblue,anchor=west] at (-4.1,3.25) {$\MH$（复维 $2n$，超凯勒）};
  % 三根纤维（环面）
  \foreach \x/\h/\lab in {-2.7/1.50/{$b_1$}, 0/2.45/{$b_2$}, 2.7/1.15/{$b_3$}}{
    \draw[thick,fill=white] (\x,\h) ellipse (0.80 and 0.30);
    \draw[thick,fill=blue!8] (\x,\h) ellipse (0.27 and 0.10);
    \draw[gray,dashed] (\x,\h-0.30)--(\x,0.60);
    \fill (\x,0.60) circle (1.2pt);
    \node[above right=2pt,font=\small] at (\x,0.60) {\lab};
  }
  % 基空间
  \draw[thick] (-4.6,0.60)--(-0.8,0.60);
  \draw[thick] (0.8,0.60)--(4.6,0.60);
  \node[font=\small,anchor=north west] at (-4.6,0.40) {$B=\bigoplus_i H^0(C,K_C^{d_i})$};
  \node[font=\small,anchor=north east] at (4.6,0.40)
    {$\omega_J|_{\mathcal F}=\omega_K|_{\mathcal F}=0$};
  % pi 箭头
  \draw[->,thick] (3.72,1.00)--(3.72,0.64);
  \node[right,font=\small] at (3.72,0.82) {$\pi$};
\end{tikzpicture}
\caption{Hitchin 纤维化。基 $B$ 是复向量空间，复维 $(g-1)\dim G$；一般纤维
$\mathcal F_b$ 是复环面（画成环面），复维与基相同：$\dim_\CC\mathcal F=(g-1)\dim G$。
$\pi$ 对复结构 $I$ 全纯，纤维在 $\omega_J,\omega_K$ 下 Lagrangian；合起来即「特殊
Lagrangian 环面纤维化」，正是 SYZ 镜像对称要求的几何。}
\label{fig:hitchin-fibration}
\end{figure}

\subsubsection*{本节自检}
\begin{enumerate}
\item 为什么 $\dim_\CC B=\tfrac12\dim_\CC\MH$ 是「可积」的标志？如果基的维数不是一半，
换句话说什么会坏掉？
\item $H_{P,\alpha}$ 中 $\alpha\in H^1(C,K_C^{1-n})$ 的次数为什么必须是这样？如果取
$\alpha\in H^0$ 或 $H^1(C,\mathcal O)$，被积式还是 $(1,1)$-型吗？
\item 纤维的 Lagrangian 性用到了哪两个事实？为什么它对 SYZ 是必需的？
\end{enumerate}

% ======================================================================
\subsection{特殊化：$\Psi=\infty$ 给 B-模型、$\Psi=0$ 给 A-模型（§5.3，pp.~66--70）}
\label{sec:fh-spec}

\subsubsection*{动机：四维理论在二维里到底选了哪一对复结构}

§\ref{sec:fh-hk} 说明 $\MH$ 上有一整个 $\PP^1_h\times\PP^1_h$ 的参数族二维 TFT，由一对
复结构 $(J_+,J_-)$ 标记；而 §\ref{sec:fh-from4d} 只说「靶空间是 $\MH$」，没说四维理论
在二维极限里\emph{选了哪一点}。这一步不能省：M2 说「四维 S-对偶
$\tau\mapsto-1/\tau$ 在二维变成镜像对称」，而镜像对称的精确内容依赖 $(J_+,J_-)$。
本小节把这一点算出来。

\subsubsection*{$(J_+,J_-)$ 语言回顾}

把 $\MH$ 上的三个复结构统一记成 $I_w$（\eqref{eq:Iw}）。给定一对点
$(w_+,w_-)\in\PP^1_h\times\PP^1_h$，二维模型由无破缺超对称方程刻画：\pA\anchor{src p.59}
\begin{equation}
(1-\ii J_+)\,\partial\Phi=0,\qquad (1-\ii J_-)\,\bar\partial\Phi=0 ,
\label{eq:susy-pair}
\end{equation}
其中 $\partial,\bar\partial$ 是 $\Sigma$ 上的 $(1,0)$ 与 $(0,1)$ 部分，$J_\pm=I_{w_\pm}$
作用在拉回切丛上。两类特别情形是：
\begin{itemize}
\item $J_+=J_-=:J_b$（即 $w_+=w_-$）：两条方程合并为 $(1-\ii J_b)\dd\Phi=0$，
  在欧氏 $\Sigma$ 上迫使 $\dd\Phi=0$，即 $\Phi$ 常映射。称
  \term{B-model}{B-模型}（在复结构 $J_b$ 中）。
\item $J_-=-J_+$（即 $w_-=-1/w_+$，用 \eqref{eq:opposite}）：两条方程互为复共轭，
  任一都断言 $\Phi$ 在 $J_+$ 下全纯。称 \term{A-model}{A-模型}（准确地说，以
  $\omega_{J_+}$ 为辛形式的 A-模型）。
\end{itemize}
一般情形 $J_-\ne\pm J_+$ 下模型是广义 B-模型（原文 §5.2 用广义复几何约化为 A-模型，
本文件不展开）。于是判断「四维理论给了什么」就归结为计算 $(w_+,w_-)$。\pA\anchor{src p.59}

\subsubsection*{参数：$\tau$、$t$ 与 canonical parameter $\Psi$}

四维理论有两个参数：规范耦合 $\tau=\theta/2\pi+4\pi\ii/g^2$，以及扭化族参数 $t$
（来自 $(3.13)$，$t$ 实对应 $F'-t^{-1}D\varphi=0$ 的超对称条件）。二维模型却只依赖
一个组合的\emph{规范参数}：\pA\anchor{src p.66}
\begin{equation}
\Psi=\frac{\tau+\bar\tau}{2}+\frac{\tau-\bar\tau}{2}\cdot
\frac{t-t^{-1}}{t+t^{-1}} .
\label{eq:Psi}
\end{equation}
为什么？$t$ 与 $\tau$ 在二维模型中不独立：改变 $B$-场（由 $\operatorname{Re}\tau$ 给出）
与改变复结构（由 $t$ 给出）在广义复几何里是同一回事（原文 §5.2 的结论：$B$ 的
$(2,0)$ 分量与复结构空间中的移动等价）。$\Psi$ 就是 $\PP^1$ 上的坐标，把两种移动
合并成一个；模掉 $t\mapsto-t$ 与 $U(1)$ 作用后，理论只依赖 $\Psi$。
$\Psi$ 是这个族里的\term{canonical parameter}{规范参数}。\pA\anchor{src p.66}

\eqref{eq:Psi} 立刻给出两个特殊点：
\begin{itemize}
\item $t=\pm\ii$ 时 $(t-t^{-1})/(t+t^{-1})=\infty$，故 $\Psi=\infty$，\emph{与 $\tau$
  无关}（弱耦合极限也一样）；
\item $t=\pm1$ 时该因子为 $0$，故 $\Psi=\operatorname{Re}\tau$；特别地
  $\Psi=0$ 当且仅当 $t=\pm1$ 且 $\theta=0$（$\operatorname{Re}\tau=0$）。
\end{itemize}
这两个点正是 S-对偶（作为 $\Psi\mapsto-1/(n_g\Psi)$ 的分式线性变换）的不动点交换对。
$|t|=1$ 时 $\Psi$ 实；$|t|\ne1$ 时 $\Psi$ 可以取到整个上半或下半平面。\pA\anchor{src p.35}

\subsubsection*{$\Psi=\infty$：$J$ 中的 B-模型}

取 $t=\ii$。由 $(5.13)$，$w_+=-t=-\ii$，$w_-=t^{-1}=-\ii$，即
$(w_+,w_-)=(-\ii,-\ii)$。而由 \eqref{eq:Iw}，$w=-\ii$ 给出 $I_w=J$。因此在
$\Psi=\infty$ 处，二维模型是 $J$ 中的 B-模型：
\begin{equation}
\text{四维扭化理论}\big|_{\Psi=\infty}\ \cong\ B_J\bigl(\MH(G,C)\bigr).
\label{eq:BJ}
\end{equation}
原文的措辞是：「To get $\Psi=\infty$, we take $t=\ii$ and thus $(w_+,w_-)=(-\ii,-\ii)$.
The resulting model is the B-model in complex structure $J$.」\pA\anchor{src p.70}

物理上这很自然：$t=\pm\ii$ 是扭化族的「B-模型极限」，B-模型的无质量场是常映射，
对应四维理论里 $F'=0$ 的解在 $C$ 上的拉回；而 $J$ 正是使平坦 $G_{\CC}$-联络的复结构
（§\ref{sec:fh-hk}），常映射的靶空间点就是 $C$ 上平坦 $G_{\CC}$-丛。两者吻合。

\subsubsection*{$\Psi=0$：$K$ 中的 A-模型}

取 $t=1$（$t=-1$ 同理）。由 $(5.13)$，$w_+=-1$，$w_-=1$。注意
$w_-=-1/w_+$（因为 $-1/(-1)=1$），所以由 \eqref{eq:opposite}，
$I_{w_-}=-I_{w_+}$：这是一对相反复结构，模型是 A-模型。具体地 $w=\pm1$ 给出
$\pm K$，所以
\begin{equation}
\text{四维扭化理论}\big|_{\Psi=0}\ \cong\ A_K\bigl(\MH(G,C)\bigr).
\label{eq:AK}
\end{equation}
原文对 $t=0,\infty$ 的情形另有说明：那时得到 A-模型在 $\pm I$；而 $t=\pm1$ 给出
A-模型在 $\pm K$。这两族在 $S$-对偶下的行为不同，是 §\ref{sec:fh-s} 中 Table 2
的来源。\pA\anchor{src p.70}

\subsubsection*{S-对偶的作用：$B_J(G)\ \longrightarrow\ A_K(\Ldual)$}

现在到了本节的核心。$S$-对偶在四维是 $\tau\mapsto-1/(n_g\tau)$；在二维它必须同时
作用在 $t$ 上，因为扭化参数不是 S-对偶不变的（原文 (3.13)、(3.55)--(3.56)）：
\begin{equation}
S:\quad \tau\longmapsto-\frac{1}{n_g\tau},\qquad
t\longmapsto-t\,\frac{\tau}{|\tau|} .
\label{eq:S-on-t}
\end{equation}
在 $\PP^1_h$ 上，由 $(5.22)$，$S$ 作用为
\begin{equation}
w_+\longmapsto w_+\,\frac{c\tau+d}{|c\tau+d|},\qquad
w_-\longmapsto w_-\,\frac{|c\tau+d|}{c\tau+d},
\label{eq:S-on-w}
\end{equation}
其中 $S=\begin{pmatrix}a&b\\c&d\end{pmatrix}$。在 $\Psi$ 上，$(3.53)$--$(3.54)$ 给出
\begin{equation}
\Psi\longmapsto-\frac{1}{n_g\Psi} .
\label{eq:S-on-Psi}
\end{equation}
注意 \eqref{eq:S-on-Psi} 是分式线性变换，把 $\Psi=0$ 与 $\Psi=\infty$ 互换。

取最简单的例子：$\tau$ 纯虚（$\theta=0$），并取 $\Psi=\infty$ 的起点 $t=\ii$。
此时 $\tau/|\tau|=\ii$，所以
\begin{equation}
t'=-t\,\frac{\tau}{|\tau|}=-\ii\cdot\ii=1,\qquad\text{即}\qquad
(w_+,w_-)=(-t,t^{-1})\longmapsto(-1,1).
\label{eq:S-derivation}
\end{equation}
而 $(w_+,w_-)=(-1,1)$ 正是 $\Psi=0$ 的 A-模型在 $K$ 中。于是\pA\anchor{src p.70}
\begin{equation}
S:\qquad B_J\bigl(\MH(G,C)\bigr)\ \cong\ A_K\bigl(\MH(\Ldual,C)\bigr).
\label{eq:S-BtoA}
\end{equation}
这条式子是 M2 的第一次精确表述。它说的是：\emph{四维 S-对偶把 $J$ 中的 B-模型变成
$K$ 中的 A-模型}。注意两组复结构不同（$J$ 与 $K$）——这正是「镜像对称」的雏形：
两个理论在同一个（或对偶的）超凯勒流形上，却选择了不同的复结构。

\subsubsection*{从 $B_J\to A_K$ 到镜像对}

\eqref{eq:S-BtoA} 的左右两边复结构不同，还不是标准形式的镜像对称（标准形式要求
$B_J$ 与 $A_J$，即同一个复结构）。原文用一步补足：在 $S$ 之后再复合理论的自同构
\begin{equation}
\mathcal I:\ \varphi\longmapsto\ii\varphi ,
\label{eq:phi-i-phi}
\end{equation}
它把复结构 $K$ 旋转回 $J$（在 $\PP^1_h$ 上绕 $I$ 轴转 $90^\circ$）。这个自同构不是
四维族 $\PP^1_g$ 的对称（它把理论移出四维紧化族），但它保持二维模型的超对称代数
（作用在 $\varphi$ 的复结构上）。复合后得到
\begin{equation}
B_J\bigl(\MH(G,C)\bigr)\ \cong\ A_J\bigl(\MH(\Ldual,C)\bigr).
\label{eq:mirror-pair}
\end{equation}
原文的结论是：\pA\anchor{src p.70}
\begin{origquote}
“We get a mirror symmetry that maps the B-model in complex structure $J$ to the A-model in complex structure $J$. So $\MH(G,C)$ and $\MH({}^LG,C)$ are a mirror pair in complex structure $J$.”
\end{origquote}
当 $G=\Ldual$（例如 $G=G_2,F_4,E_8$，或自对偶的 $SU(N)$ 在特殊情形）时两个流形
是同一个，\eqref{eq:mirror-pair} 就是 $\MH$ 上\emph{到自身的}镜像对称。一般 $G$ 时
两个流形不同，但它们的 Hitchin 纤维化有相同的基（下一小节）。

\subsubsection*{Table 2：S-对偶在六个特殊点上的作用}

把上面两步合起来，就得到原文 Table 2（p.71）的内容。表中每一行是一个二维模型，
第二列是它在 S-对偶下的像：\pA\anchor{src p.71}
\begin{center}
\begin{tabular}{@{}ll@{}}
\toprule
起点 & 像 \\
\midrule
$I_B$ & $I_B$ \\
$I_A$ & $I_A$ \\
$J_B$ & $K_A$ \\
$J_A$ & $K_B$ \\
$K_B$ & $J_A$ \\
$K_A$ & $J_B$ \\
\bottomrule
\end{tabular}
\end{center}
注意 $I_B,I_A$ 是 S-对偶的\emph{不动点}：$I_w$ 在 $w=0$（$I$）与 $w=\infty$（$-I$）
处被 $U(1)_R$ 固定（§\ref{sec:fh-hk}），而 $S$ 作用在这个 $U(1)_R$ 上。所以
「$B$ 换成 $A$」只发生在 $J,K$ 那一对——$I$ 对 S-对偶是透明的。这解释了为什么
§5.4 的论证可以在 $I$ 里做（$I_B$ 不动），而结论却落在 $J_B$ 与 $K_A$ 上。
\pA\anchor{src p.71}

\begin{remark}[镜像对称的一个独立检验]
当 $G=\Ldual$ 时，任意把 $B_J$ 映到 $A_J$ 的镜像对称 $m$ 必保持上同调环的 cup
product（因 $\alpha\cdot\beta$ 与 $\alpha\cdot m^*\beta$ 都非零），却把 $(p,q)$ 型送到
$(n-q,n-p)$ 型（$n=\dim_\CC\MH$）。于是 B-模型的 $(p,q)$-算符数必须等于 A-模型的，
这是 \eqref{eq:mirror-pair} 非平凡的旁证。\pB\anchor{src p.71}
\end{remark}

\subsubsection*{图：$(\tau,\Psi)$ 平面与 S-对偶}

\begin{figure}[htbp]
\centering
\begin{tikzpicture}[scale=0.85,>={Stealth[length=2mm]}]
  % 左：τ 上半平面
  \begin{scope}
    \draw[thick] (-2.2,0)--(2.2,0);
    \node[below,font=\small] at (2.15,-0.02) {$\operatorname{Re}\tau$};
    \node[below left=0pt,font=\small] at (0,0) {$0$};
    \node[font=\small,align=center] at (1.35,1.52) {$\tau$ 上半平面};
    % 基本区域
    \fill[blue!8] (-0.5,0)--(0.5,0)--(0.5,1.4)--(-0.5,1.53)--(-0.5,0);
    \draw[blue!60] (-0.5,0)--(-0.5,1.53) (0.5,0)--(0.5,1.4);
    \draw[blue!60] (-0.5,1.53) arc (150:30:1.0 and 1.53);
    \fill (0,1.0) circle (1.2pt) node[right=2pt,font=\scriptsize] {$\ii$};
    % S 与 T
    \draw[->,thick,kwblue] (-1.75,0.25) to[bend left=25] node[above,font=\scriptsize] {$S$} (-0.25,1.75);
    \draw[->,thick,kwgreen] (1.85,0.35)--(0.85,0.35) node[midway,above,font=\scriptsize] {$T$};
  \end{scope}
  % 右：Ψ 球面
  \begin{scope}[xshift=6.2cm]
    \draw[thick] (0,0) circle (1.45);
    \draw[gray!70] (-1.45,0) arc (180:360:1.45 and 0.42);
    \draw[gray!70,dashed] (1.45,0) arc (0:180:1.45 and 0.42);
    \fill (0,1.45) circle (1.4pt) node[above=1pt,font=\small] {$\Psi=\infty$};
    \fill (0,-1.45) circle (1.4pt) node[below=1pt,font=\small] {$\Psi=0$};
    % 实轴
    \draw[gray,dashed] (-1.45,0)--(1.45,0);
    \node[font=\scriptsize] at (1.02,0.16) {$|t|=1$};
    \draw[->,thick,kwblue] (-0.75,1.05) to[bend right=40] node[above=1pt,font=\scriptsize] {$S$} (0.75,1.05);
    \node[font=\scriptsize] at (-0.75,0.52) {$|t|>1$};
    \node[font=\scriptsize] at (-0.75,-0.52) {$|t|<1$};
  \end{scope}
\end{tikzpicture}
\caption{左：$S$ 与 $T$ 在 $\tau$ 上半平面上的作用（基本区域用阴影画出，$S$ 把
$\tau=\ii$ 附近送向原点方向）。右：canonical parameter $\Psi\in\PP^1$。
$S$ 在 $\Psi$ 上是 $\Psi\mapsto-1/(n_g\Psi)$，把 $\Psi=0$（南极点，A-模型在 $K$ 中）
与 $\Psi=\infty$（北极点，B-模型在 $J$ 中）互换。赤道 $|t|=1$ 是 $\Psi$ 取实值的
轨迹（$\Psi=\operatorname{Re}\tau$）；上、下半球分别对应 $|t|>1$ 与 $|t|<1$。}
\label{fig:tau-psi}
\end{figure}

\paragraph{排除错误直觉之三：S-对偶不是把 $\Psi\to-1/\Psi$ 当作「另一个理论的参数」，
而是同一个二维理论族的坐标变换。}
$\Psi$ 是 $\PP^1$ 上的坐标，不是某个流形的模参数。S-对偶把 $\Psi$ 变成
$-1/(n_g\Psi)$，等价于把族中 $\Psi=0$ 的点与 $\Psi=\infty$ 的点对换——这正是
「B-模型与 A-模型互换」的精确含义。所以「四维 S-对偶在二维是镜像对称」不是一句
比喻，而是一个可计算的坐标变换。\pA\anchor{src p.67}

\subsubsection*{本节自检}
\begin{enumerate}
\item \eqref{eq:Psi} 中为什么 $\Psi$ 只依赖 $t$ 的比 $(t-t^{-1})/(t+t^{-1})$ 与 $\tau$？
如果 $t$ 与 $\tau$ 独立，二维模型会依赖两个参数，那会与 §\ref{sec:fh-from4d} 的哪个结论冲突？
\item 用 \eqref{eq:Iw} 验证 $w=-\ii$ 给 $J$、$w=\pm1$ 给 $\pm K$，从而核对
\eqref{eq:BJ} 与 \eqref{eq:AK}。
\item 为什么 $I_B,I_A$ 在 Table 2 中是 S-对偶的不动点？这与 §\ref{sec:fh-hk} 中
$U(1)_R$ 固定 $w=0,\infty$ 有什么关系？
\end{enumerate}

% ======================================================================
\subsection{S-对偶在 Hitchin 纤维化上的作用：基上经典、纤维上 T-对偶（§5.4--5.5，pp.~71--77）}
\label{sec:fh-s}

\subsubsection*{主命题}

上一小节给出了 M2 的\emph{代数}形式：$B_J(G)\cong A_J(\Ldual)$。但镜像对称是一个
关于\emph{范畴}的陈述（B-模型与 A-模型的 D-膜范畴等价），仅凭一个代数等式不足以
相信它。§5.4--5.5 给了两个独立的论证：一个用 B-模型的 chiral ring（基上），
一个用膜（纤维上）。合起来就是 SYZ 图像。原文把主命题写在 §5.4 开头：\pA\anchor{src p.71}
\begin{origquote}
“Our aim here and in section 5.5 is to show that S-duality acts classically on the base of the Hitchin fibration $\pi:\MH\to B$, while acting by T-duality on the fibers.” \anchor{src p.71}
\end{origquote}
这句话是全节的目标，也是本文件标题「S-对偶为什么就是镜像对称」的答案。

\subsubsection*{基上：chiral ring 论证}

先看基。B-模型在复结构 $I$ 中的 chiral ring 是
$H^q(X,\wedge^pT_X)$（双分次），其中 $p=q=0$ 的子环是 $X$ 上全纯函数的环。对
$X=\MH$ 在复结构 $I$ 下，这正是 $\MH$ 的全纯函数环。而 Hitchin 纤维化的纤维是
\emph{紧}复子流形（§\ref{sec:fh-fib}），所以 $\MH$ 上的全纯函数在每条纤维上为常数，
必是基 $B$ 上全纯函数的拉回：\pA\anchor{src p.72}
\begin{equation}
\Gamma\bigl(\MH,\mathcal O_{\MH}\bigr)\ \cong\ \Gamma\bigl(B,\mathcal O_B\bigr).
\label{eq:chiral-base}
\end{equation}
于是 B-模型的局域算符环只看见基 $B$，看不见纤维。若 S-对偶是 B-模型的范畴等价
（从 $G$ 到 $\Ldual$），则它诱导
\begin{equation}
\Gamma(B,\mathcal O_B)\ \cong\ \Gamma(\Ldual B,\mathcal O_{\Ldual B})
\label{eq:ring-iso}
\end{equation}
一个分次环同构。这就证明了「S-对偶在基上经典地作用」：它给出基之间的一个映射
$\Xi:B\to\Ldual B$。\pB\anchor{src p.72}

$\Xi$ 具体是什么？原文的论证走四维：四维理论的半 BPS 算符 $D_n$ 对应
$\mathrm{Sym}^n\mathbf 6$（$\mathbf 6=\wedge^2\mathbf 4$ 是 $SU(4)_R$ 的六维表示），
其最高权分量是 $G$-不变多项式 $P(\varphi)$。扭化后 $D_n$ 的共形维数为零、OPE 无奇异，
因此构成一个 $\CC^*$-分次环 $D'$。S-对偶把 $D_n(G)$ 送到 $D_n(\Ldual)$，故给出
$D'$ 的分次环同构；对单连通的 $G$ 这个作用是非平凡外自同构的补足（$G_2,F_4$ 上是
非平凡对合，一般 $G$ 由 $\Ldual$ 替换），在取好归一化后可以取为恒等。于是
$\Xi$ 由「$S$ 在基的生成元 $P_i(\varphi)$ 上作用」唯一确定：它是一个
\emph{仿射（线性）变换}，由 $\CC^*$ 的伸缩 $P_i\mapsto\lambda^{d_i}P_i$ 与
$\lambda\mapsto\lambda/\Psi$ 复合而成。\pB\anchor{src p.73}

基 $B$ 是向量空间，其坐标是 Casimir 不变量的值；S-对偶只能按次数重新标度它们
（$P_i$ 是 $d_i$ 次齐次），不能做非线性变换——否则会破坏 $B$ 的线性结构，而线性结构
正是 $U(1)$ 作用与可积系统 Hamiltonians 的来源。所以「基上的 S-对偶 $=$ 仿射变换
$\Xi$」不是假设，是 $B$ 的线性结构逼出来的。\pB\anchor{src p.73}

\subsubsection*{纤维上：为什么必须是 T-对偶}

基上没问题了，纤维上却出了麻烦。原文的原话是：\pA\anchor{src p.72}
\begin{origquote}
“there is no analogous way to show that S-duality acts classically on the fibers of the Hitchin fibration. It hardly can, given its relation to mirror symmetry!”
\end{origquote}
推理是：在 B-模型（复结构 $I$）里局域算符都是基上的全纯函数（\eqref{eq:chiral-base}），
它们测不到纤维上的位置；纤维上的动量对应 massive 态，而不是局域算符。所以 S-对偶
对纤维的作用\emph{不可能是经典映射}。而 SYZ 图像恰好要求：镜像对称在纤维上不是经典
映射，而是 T-对偶——在每一条环面纤维上取对偶环面。本情形下（\eqref{eq:ring-iso} 已
给出 $\Xi$）我们\emph{知道}基映射是什么，所以 SYZ 图像可以精确化：\pA\anchor{src p.77}
\begin{origquote}
“This is the usual SYZ picture of mirror symmetry, except that the map $\Xi$ is usually assumed to be trivial.” \anchor{src p.77}
\end{origquote}

\subsubsection*{膜论证：点膜在 S-对偶下的像必须支撑在对偶纤维上}

上面的论证是「不可能经典」，但还没证明「一定是 T-对偶」。§5.5 用零维膜把这个缺口
补上，这是本文最漂亮的一段推理，我们完整重建如下。\pA\anchor{src p.75}

\textbf{第一步：点 = $(B,B,B)$-膜。}$\MH$ 上的一个零维膜（点 $p$）没有标量场期望值，
其边界条件就是固定 $\Phi\equiv p$。点在任何复结构下都是复子流形（维数 $0$），配上
平凡线丛，所以它是 $(B,B,B)$-型的：相对于 $I,J,K$ 都是 B-膜。

\textbf{第二步：S-对偶后类型变成 $(B,A,A)$。}查 Table 2，$(B,B,B)\mapsto(B,A,A)$。
所以在对偶理论 $\MH(\Ldual)$ 中，像膜 $B_p$ 必须相对于 $I$ 是 B-膜、相对于
$J,K$ 是 A-膜。

\textbf{第三步：支撑维数两边夹。}A-膜在其辛形式下要求支撑至少是 middle-dimensional
（因为 A-模型需要 $2\,\Re\Omega$ 沿支撑非退化）；而 $B_p$ 的模空间（形变空间）只是
一个点，维数为 $0$，所以支撑至多 middle-dimensional。两边合并：$B_p$ 的支撑
\emph{恰好}是 middle-dimensional，即复维 $(g-1)\dim G$。

\textbf{第四步：中点就是纤维。}对于任意点 $q\in\MH(\Ldual)$，若 $q$ 在 $B_p$ 的支撑
上，则 $q$ 与 $p$ 有相同的 Casimir 不变量（因为 $p$ 是固定点，$B_p$ 只能沿基的
固定值移动），所以 $B_p$ 的支撑落在 $\Ldual\pi^{-1}(\Xi(\pi(p)))$ 内。两者都是
middle-dimensional，故支撑就是那条纤维。

\textbf{第五步：纤维上的平移。}点膜 $B_p$ 上的开弦无质量算符是 $\Phi$ 的取值函数；
而 $B_p$ 的支撑是纤维，$\Phi$ 在纤维上取遍所有值，所以 $B_p$ 的模空间就是纤维
$\mathcal F$ 本身。对偶膜 $\widehat B_p$ 的模空间则是「$\mathcal F$ 上的平坦
Chan-Paton 线丛」的模空间，即对偶环面 $\widehat{\mathcal F}$。于是 S-对偶给出
同构
\begin{equation}
\mathcal F\ \cong\ \widehat{\mathcal F}:=\operatorname{Hom}\bigl(\pi_1(\mathcal F),U(1)\bigr),
\label{eq:SYZ-fiber}
\end{equation}
即「纤维与对偶纤维是 T-对偶」。原文强调 rank $1$（不是更高 rank）：若取 rank $>1$
的 Chan-Paton 丛，$\widehat B_p$ 的模空间维数会大于 $B_p$，S-对偶就不可能是同构。
\pA\anchor{src p.76}

\begin{keyidea}
\textbf{整个论证链一句话。}点膜没有复结构可选（点在任何复结构下都全纯），故必为
$(B,B,B)$；S-对偶把它变成 $(B,A,A)$，而 $(B,A,A)$ 的 A-膜性质强迫其支撑为
middle-dimensional；于是像点「被迫」长大成一条纤维。这就是镜像对称把点送到环面的
机制，也是 SYZ 从四维理论里\emph{被推导出来}（而不是被假设）的原因。
\end{keyidea}

\subsubsection*{数学依据与微妙处}

原文指出 $\mathcal F\cong\widehat{\mathcal F}$ 在文献中已对若干情形独立知道：$G$ 为
酉群时由 Hausel--Thaddeus（[21]），$G=G_2$ 由 Hitchin（[81]），一般半单 $G$ 由
Donagi--Pantev（[82]）。本节不重新证明它们，只说明它们都是同一个四维 S-对偶的不同
侧面。微妙处：$G$ 的中心与基本群在 §7 才起作用（那时 $\Xi$ 不再是恒等）；$C$ 上的
spin 结构本文不分析。\pA\anchor{src p.77}


\paragraph{排除错误直觉之四：T-对偶发生在纤维上，不在底空间上。}
SYZ 镜像对称的 T-对偶作用在环面\emph{纤维}上，基 $B$ 上的映射 $\Xi$ 是经典的
（光滑甚至仿射的）。如果把 T-对偶理解成基上的变换，就会得到「镜像对称把 $B$ 映到
$B$ 之外」的错误结论。原文的措辞很明确：「S-duality acts classically on the base …
while acting by T-duality on the fibers」。基上可测（有局域算符），纤维上不可测
（无局域算符）——这两句话是同一件事。\pA\anchor{src p.71}

\paragraph{排除错误直觉之五：S-对偶在基上不是「把一个点送到另一个点」的任意映射。}
$\Xi$ 受两个约束：（i）它是分次环同构诱导的，因此保持基的线性结构（仿射）；（ii）
它把 $P_i$ 按次数 $d_i$ 标度，故与 $U(1)$ 作用相容。这两条使 $\Xi$ 不是任意的双射，
而是「$\CC^*$ 伸缩 + 可能的 $\Psi$ 依赖」的复合。相应地，S-对偶在 $\Psi$ 上的分式
线性作用 \eqref{eq:S-on-Psi} 不是 $\Xi$ 的作用，而是族参数的作用——两者不要混淆。
\pB\anchor{src p.73}

\subsubsection*{本节自检}
\begin{enumerate}
\item 为什么「局域算符测不到纤维」会迫使 S-对偶在纤维上取 T-对偶，而不是别的
  （例如恒等）？如果 S-对偶在纤维上取恒等，与 $\mathcal F\cong\widehat{\mathcal F}$
  矛盾在哪里？
\item 点膜论证的第三步用到了「A-膜支撑至少 middle-dimensional」与「点膜支撑至多
  middle-dimensional」。这两个「维数界」各自的来源是什么？
\end{enumerate}

% ======================================================================
\subsection{$\MH$ 上的膜分类与本文对镜像对称概念的贡献（§5.6，pp.~77--79）}
\label{sec:fh-brane}

\subsubsection*{类型 $(B,B,B)$、$(B,A,A)$、$(A,B,A)$、$(A,A,B)$}

$\MH$ 上有三个复结构 $I,J,K$，因此一张膜相对于每一个复结构都可以问「它是 A-型还是
B-型」。把三个答案写成三元组：$(B,B,B)$ 表示相对于 $I,J,K$ 都是 B-膜；$(B,A,A)$
表示对 $I$ 是 B-膜、对 $J,K$ 是 A-膜；等等。原文的清单是：\pA\anchor{src p.77}
\begin{itemize}
\item $(B,B,B)$：B-膜对每个复结构。点膜是例子；更一般地，space-filling 的
  全纯向量丛也给出这一类。
\item $(B,A,A)$：I-全纯的复子流形 $N$，且对全纯辛形式 $\Omega_I$ 是 Lagrangian 的
  （从而也对 $\omega_J,\omega_K$ Lagrangian）。例子：Hitchin 纤维
  $\mathcal F_b$；固定 $E$ 让 $\varphi$ 变化的 $T^*\mathcal M$ 纤维（此时
  $\Omega_I=\int\Tr\delta\varphi_z\wedge\delta A_z$ 在固定 $A$ 时为 $0$）。
\item $(A,B,A)$ 与 $(A,A,B)$：对 $J$（或 $K$）的复 Lagrangian 子流形。例子：固定
  $A_z=A_z+\ii\varphi_z$ 让 $A_{\bar z}$ 变化，即指定了全纯结构的平坦 $G_{\CC}$-丛。
\end{itemize}
注意 $(B,A,A)$ 的例子有一个统一的形式：$N$ 是复结构 $I$ 下的复子流形，且在
$\Omega_I$ 下 Lagrangian。这时它自动对 $\omega_J,\omega_K$ 也 Lagrangian
（因为 $\Omega_I=\omega_J+\ii\omega_K$ 在 $N$ 上为零意味着两个实部都为零）。
\pA\anchor{src p.78}

\subsubsection*{哪些类型不存在}

不是所有三元组都能实现。原文脚注 24 指出：\pA\anchor{src p.78}
\begin{origquote}
“there are no branes of type $(B,B,A)$: a subvariety holomorphic in complex structures $I$ and $J$ is automatically holomorphic in complex structure $K$.”
\end{origquote}
理由很干净：若 $N$ 同时是 $I$-全纯与 $J$-全纯的，由四元数关系 $K=IJ$，$N$ 自动
$K$-全纯；但「对 $K$ 是 A-膜」要求 $N$ 对 $\omega_K$ 迷向（Lagrangian），而一个
$K$-全纯子流形上 $\omega_K$ 的限制是 Kähler 形式，只有在 $N$ 为离散点时才为零——
所以非平凡的 $(B,B,A)$ 不存在。$(A,A,A)$ 同理：A-膜要求支撑 middle-dimensional，
三对相互矛盾的迷向性无法同时满足。\pB\anchor{src p.78}

\subsubsection*{S-对偶如何置换类型}

结合 Table 2，S-对偶对膜类型的作用是：\pA\anchor{src p.78}
\begin{equation}
(B,B,B)\longleftrightarrow(B,A,A),\qquad
(A,B,A)\longmapsto(A,B,A),\qquad
(A,A,B)\longmapsto(A,A,B).
\label{eq:brane-permutation}
\end{equation}
第一行正是点膜论证用的那一条：S-对偶把点膜变成支撑在纤维上的 $(B,A,A)$ 膜。第二、三行
说 $(A,B,A)$ 与 $(A,A,B)$ 是 S-对偶不变的\emph{类型}（具体的膜被送到另一张具体的膜，
类型不变）；这与 \eqref{eq:S-BtoA} 中 $J\leftrightarrow K$ 的交换一致——「B」所指的
复结构在两侧按 $\Xi$ 对应，故不冲突。\pB\anchor{src p.78}

\subsubsection*{本文对镜像对称概念的贡献}

现在可以回答本文件的中心问题：\emph{为什么说 S-对偶就是镜像对称？}

\textbf{第一，镜像对称在这里不是两个 Calabi--Yau 流形之间的等价。}标准表述是：
两个 CY 流形 $X,\hat X$ 的 A-模型与 B-模型等价。本文的情形不同：$X=\MH(G,C)$ 与
$\hat X=\MH(\Ldual,C)$ 是\emph{超凯勒}流形（非紧、带奇点），而且当 $G=\Ldual$
时它们\emph{是同一个流形}。镜像对称变成 $\MH$ 上两个不同复结构对应的 A/B 模型之间的
等价：$B_J\leftrightarrow A_J$（\eqref{eq:mirror-pair}）。所以它是\emph{同一个几何对象
上的两种量子场论}的等价，而不是两个几何对象之间的对应——这正是本文的贡献之一。
\pA\anchor{src p.70}

\textbf{第二，SYZ 纤维化是被推导出来的，不是被假设的。}标准 SYZ 假设镜像流形带特殊
Lagrangian 环面纤维化、且基上映射平凡。本文里 Hitchin 纤维化\emph{自动}提供这个结构
（§\ref{sec:fh-fib}），基映射 $\Xi$ 由 chiral ring 论证\emph{算出}
（\eqref{eq:ring-iso}），纤维上的 T-对偶由点膜论证\emph{证明}
（\eqref{eq:SYZ-fiber}）。原文说「except that the map $\Xi$ is usually assumed to be
trivial」（p.77），意思正是：一般 SYZ 中 $\Xi$ 平凡是假设，这里它由四维理论给出。
\pA\anchor{src p.77}

\textbf{第三，A/B 的「复结构选择」获得了四维解释。}在 $\MH$ 上选 $J$ 还是 $K$ 不是
任意的：$J$ 是 $\Psi=\infty$ 的 B-模型，$K$ 是 $\Psi=0$ 的 A-模型
（§\ref{sec:fh-spec}），而 S-对偶作用 $\Psi\mapsto-1/(n_g\Psi)$ 正好交换这两点。所以
「镜像对称交换 A 与 B」在四维就是一个关于 $\tau,t$ 的坐标变换——这是解释，不是类比。
\pA\anchor{src p.70}

\paragraph{排除错误直觉之六：B-模型不是「裸的拓扑 B-模型」的原样搬用。}
二维 B-模型的标准定义只需要靶空间是一个复流形；而在扭化理论里，B-模型还带着
B-场，B-场的 $(2,0)+(0,2)$ 分量与复结构一起构成\emph{广义复结构}（原文 §5.2）。
只有当 B-场取特殊值（或等价地，$\Psi$ 取特殊值）时，广义复结构才退化为普通复结构，
模型才是「纯」B-模型。所以「$\MH$ 上的 B-模型」有两层意思：（i）复结构 $I$（或 $J$）
中的纯 B-模型；（ii）带 B-场的广义 B-模型。§5.3 的特殊化 $\Psi=\infty$ 与 $\Psi=0$
正是从（ii）退化到（i）的两个点。混淆这两层会把 Table 2 读错。\pB\anchor{src p.69}

\paragraph{排除错误直觉之七：镜像对称不是「两个 CY 流形之间的等价」这么简单。}
如第一点所述，本情形中两个流形可以相同（$G=\Ldual$），镜像对称是同一个流形上
不同复结构对应的模型之间的等价；更准确地说，它是\emph{二维 TFT 族}（由 $(J_+,J_-)$
标记）内部的一个对称性，而四维 S-对偶是这个族上的一个具体变换；读成「$X$ 与
$\hat X$ 的对应」会漏掉「同一流形、不同复结构」这一层。\pA\anchor{src p.77}

\subsubsection*{本节自检}
\begin{enumerate}
\item 为什么 $(B,A,A)$ 膜的支撑自动对 $\omega_J,\omega_K$ Lagrangian？用
  $\Omega_I=\omega_J+\ii\omega_K$ 写出来。
\item 脚注 24 的论证里，哪一步用到了「$K=IJ$」？如果 $I,J,K$ 不满足四元数关系，
  结论会怎样？
\item 本文说「SYZ 在这里是定理而不是假设」，具体指哪两个论证？它们各自证明了什么？
\end{enumerate}

% ======================================================================
\subsection*{本文件在全书中的位置}
\noindent\textbf{S3}（\S\ref{sec:fh-from4d}--\ref{sec:fh-fib}）：四维理论在 $\Sigma\times C$ 上紧化，靶空间 $\MH$；$\MH$ 是超凯勒商、带 $\PP^1_h$ 族复结构与 Hitchin 纤维化，产出「$\MH$ 上的二维 TFT 族」。\textbf{S4}（\S\ref{sec:fh-spec}--\ref{sec:fh-brane}）：$\Psi=\infty$ 给 $B_J$、$\Psi=0$ 给 $A_K$，S-对偶交换它们，即 $B_J(G)\cong A_J(\Ldual)$（这就是 M2），产出「$\MH$ 上的 A/B 模型范畴及其等价」。\textbf{S5 之后}：线算子作为 $\MH$ 上的膜算子来读，类型正是 \S\ref{sec:fh-brane} 的 $(B,A,A)$ 等；eigenbrane 是 Hitchin 纤维上的 $(B,A,A)$ 膜，Hecke 特征层由固定 $A_z$ 的 $(A,B,A)$ 膜构造，D-模的配对落到 \S\ref{sec:fh-s} 的 SYZ 对偶上——没有 \S\ref{sec:fh-hk} 的三复结构与 \S\ref{sec:fh-fib} 的纤维化，这些步骤都无法陈述。

若只想记住一件事：\emph{$\MH$ 上的镜像对称 $=$ 同一个超凯勒流形上 $J$ 与 $K$ 两个复结构之间的 A/B 交换 $=$ 四维 S-对偶在二维的像}；而这一切之所以可能，是因为 $\MH$ 是超凯勒的，而它超凯勒是因为它是四维 Hitchin 方程解的超凯勒商。
